\resizebox{\linewidth}{!}{
\begin{tikzpicture}[
  box/.style={rectangle, rounded corners=4pt, text width=3.4cm, minimum height=1.6cm, align=center, inner sep=6pt, font=\small},
  atk/.style={box, fill=red!10, draw=red!50!black},
  sys/.style={box, fill=orange!12, draw=orange!60!black},
  cmp/.style={box, fill=blue!8, draw=blue!45!black},
  grplbl/.style={font=\scriptsize\bfseries, text=gray!60!black}
]

% 背景分組（上排）
\fill[red!5, rounded corners=6pt] (-0.25, 0.85) rectangle (11.15, 3.3);
% 背景分組（下排左）
\fill[orange!6, rounded corners=6pt] (-0.25, -2.0) rectangle (7.55, 0.35);
% 背景分組（下排右）
\fill[blue!5, rounded corners=6pt]  (7.45, -2.0) rectangle (11.15, 0.35);

% 分組標籤
\node[grplbl] at (5.45, 3.1)  {外部攻擊威脅};
\node[grplbl] at (3.55, 0.15) {設計／使用風險};
\node[grplbl] at (9.3,  0.15) {法規風險};

% 上排：外部攻擊
\node[atk] at (1.75, 2.05) {
  \textbf{提示詞注入}\\[2pt]
  外部內容操控 Agent\\[-1pt]
  {\scriptsize\color{red!70!black}→ 情境一}
};
\node[atk] at (5.45, 2.05) {
  \textbf{供應鏈攻擊}\\[2pt]
  惡意套件植入後門\\[-1pt]
  {\scriptsize\color{red!70!black}→ 情境二}
};
\node[atk] at (9.15, 2.05) {
  \textbf{上下文操控}\\[2pt]
  安全守則遭覆蓋\\[-1pt]
  {\scriptsize\color{red!70!black}→ 情境三}
};

% 下排：系統性
\node[sys] at (1.75, -1.0) {
  \textbf{過度代理}\\[2pt]
  超越授權自主執行\\[-1pt]
  {\scriptsize\color{orange!80!black}→ 情境四}
};
\node[sys] at (5.45, -1.0) {
  \textbf{資料外洩}\\[2pt]
  憑證個資洩漏外部\\[-1pt]
  {\scriptsize\color{orange!80!black}→ 情境一、二}
};
\node[cmp] at (9.15, -1.0) {
  \textbf{合規風險}\\[2pt]
  個資法規法律責任\\[-1pt]
  {\scriptsize\color{blue!70!black}→ 全部情境}
};

\end{tikzpicture}
}
